%% The line mask: the list, the frame, and what it costs.
%% Owner: r11-mask.  Carries \label{app:linemask} (cited by 04_method.tex and
%% app_C_spatialnorm.tex).  \label{sec:linecost} is GONE: it sat on an
%% \emph{} run-in heading, so it anchored on the appendix itself and was a
%% second name for \label{app:linemask} -- every \ref to it printed C and
%% nothing warned.
\section{The line mask}
\label{app:linemask}

A crossing that falls on a known molecular transition has an astrophysical
explanation to hand, and the mask exists to say which crossings those are. It
is a label and not a verdict: \emph{every} crossing goes on to the spatial and
recurrence tests whether or not it carries the label, so a transmitter that
happens to lie near a transition is not excluded by construction.

\emph{The list.} The mask holds every transition of the \FrNSpec{} species
these fields are observed to emit --- the three CO isotopologues, HCN,
HCO$^{+}$, CS, CN, SiO, H$_2$CO, SO, [C\,\textsc{i}] and hydrogen
recombination --- that falls inside one of the \McNIsland{} searched frequency
intervals with an upper-state energy below \McEuMax\,K. That rule is the whole
of it, and it admits \FrNTrans{} transitions (Table~\ref{tab:masklines}). No
strength cut is applied: the Einstein coefficients of the CO ladder are
$10^{-7}$--$10^{-6}$\,s$^{-1}$, so any threshold admitting the strong lines of
HCN or SO would discard the brightest lines in these fields. The rarer
isotopologues, $^{34}$SO and H$^{13}$CN and the rest, are not among the named
species and would add \McNRare{} transitions. The list is a Splatalogue query
of CDMS and JPL \citep{Pickett1998,CDMS2005}, frozen with the data release, so
that the mask is reproducible as a file and not as a description.

\emph{The frame is the star's own, and the sign is not the usual one.} A
crossing's frequency is carried from the topocentric value to the barycentre
using its own block's pointing and mid-time, and then to the stellar rest
frame using the star's systemic velocity. The barycentric term runs from
\FrBaryLo{} to \FrBaryHi\,km\,s$^{-1}$ across the survey, so a tolerance
applied in the observed frame would spend \FrBarySwing\,km\,s$^{-1}$ of its
budget on the Earth's motion. The \FrEpochNEpoch{} \FrEpochStar{} epochs show
the transform working in the only way that can be checked:
\FrEpochTopoKms\,km\,s$^{-1}$ apart as observed, they agree to
\FrEpochStelKms\,km\,s$^{-1}$ in the star's frame. Throughout,
$\Delta v_\star = c\,(\nu_\star-\nu_0)/\nu_0$, so a positive offset means the
crossing lies \emph{above} the transition in frequency --- the opposite sign
to the velocity convention of a spectral line.

\emph{The half-width, and what it costs.}
A one-channel test is too tight for circumstellar gas, whose emission is
shifted by many channels: the tolerance follows from Keplerian motion where
these molecules survive, $\lesssim$13\,km\,s$^{-1}$ beyond 10\,au, and is
$\pm\MaskHalfKms$\,km\,s$^{-1}$. Between $\pm\FrLadderLo$ and
$\pm\FrLadderHi$\,km\,s$^{-1}$ the mask labels \FrAttrWTwenty{} to
\FrAttrWHundred{} of the \NCross{} crossings and removes \FrMaskPctLo{} to
\FrMaskPctHi{} per cent of the searched union, while the \FrNRankPass{}
crossings that outrank all \NCtrl{} controls are the same two at every width
and neither recurs. At the adopted width the cost is \FrMaskLostGHz\,GHz, or
\FrMaskPct{} per cent of the \FrMaskUnionGHz\,GHz unique-frequency union,
concentrated in the line-tuned fine-channel windows; the chain against the
half-width and the cost band by band are both tabulated in the data release.

\begin{table*}
\centering
\scriptsize
\caption{The line mask, every transition it holds. $n$ counts the hyperfine
components of one transition. Values are CDMS, except the hydrogen
recombination lines, which carry no $E_{\rm u}$ (K).}
\label{tab:masklines}
%% GENERATED by maskcat_v412.py -- do not hand-edit.
\begin{tabular}{@{}llr@{\,}rr@{}}
\hline
Species & Transition & $\nu_0$ (GHz) & $E_{\rm u}$ & $n$ \\
\hline
H & $41\alpha$ & 92.034434 & -- &  \\
SO & $4_{5}$--$4_{4}$ & 100.029640 & 39 &  \\
H$_2$CO & $6_{1,5}$--$6_{1,6}$ & 101.332991 & 88 &  \\
H & $39\alpha$ & 106.737357 & -- &  \\
SO & $2_{3}$--$1_{2}$ & 109.252220 & 21 &  \\
C$^{18}$O & $1$--$0$ & 109.782173 & 5 &  \\
CN & $1_{1/2}$--$0_{1/2}$ & 113.144157 & 5 & 4 \\
CN & $1_{3/2}$--$0_{1/2}$ & 113.490970 & 5 & 5 \\
CO & $1$--$0$ & 115.271202 & 6 &  \\
H & $38\alpha$ & 115.274399 & -- &  \\
H$_2$CO & $2_{0,2}$--$1_{0,1}$ & 145.602949 & 10 &  \\
CS & $4$--$3$ & 195.954211 & 24 &  \\
SO & $2_{2}$--$3_{2}$ & 210.202557 & 19 &  \\
H & $31\alpha$ & 210.501771 & -- &  \\
SO & $7_{8}$--$7_{7}$ & 214.357039 & 81 &  \\
SO & $5_{5}$--$4_{4}$ & 215.220653 & 44 &  \\
SiO & $5$--$4$ & 217.104919 & 31 &  \\
H$_2$CO & $3_{0,3}$--$2_{0,2}$ & 218.222192 & 21 &  \\
\hline
\end{tabular}%
\hfill
\begin{tabular}{@{}llr@{\,}rr@{}}
\hline
Species & Transition & $\nu_0$ (GHz) & $E_{\rm u}$ & $n$ \\
\hline
H$_2$CO & $3_{2,2}$--$2_{2,1}$ & 218.475632 & 68 &  \\
H$_2$CO & $3_{2,1}$--$2_{2,0}$ & 218.760066 & 68 &  \\
C$^{18}$O & $2$--$1$ & 219.560354 & 16 &  \\
SO & $6_{5}$--$5_{4}$ & 219.949442 & 35 &  \\
$^{13}$CO & $2$--$1$ & 220.398701 & 16 & 4 \\
H$_2$CO & $3_{1,2}$--$2_{1,1}$ & 225.697775 & 33 & 7 \\
CN & $2_{3/2}$--$1_{3/2}$ & 226.359871 & 16 & 7 \\
CN & $2_{3/2}$--$1_{1/2}$ & 226.659558 & 16 & 5 \\
CN & $2_{5/2}$--$1_{3/2}$ & 226.874781 & 16 & 6 \\
CN & $2_{5/2}$--$1_{1/2}$ & 227.191820 & 16 &  \\
CO & $2$--$1$ & 230.538000 & 17 &  \\
H & $30\alpha$ & 231.900928 & -- &  \\
CS & $5$--$4$ & 244.935557 & 35 &  \\
SO & $2_{3}$--$3_{2}$ & 246.404588 & 21 &  \\
C$^{18}$O & $3$--$2$ & 329.330553 & 32 &  \\
SO & $1_{2}$--$0_{1}$ & 329.385477 & 16 &  \\
$^{13}$CO & $3$--$2$ & 330.587982 & 32 & 4 \\
SO & $10_{11}$--$10_{10}$ & 336.553811 & 143 &  \\
\hline
\end{tabular}%
\hfill
\begin{tabular}{@{}llr@{\,}rr@{}}
\hline
Species & Transition & $\nu_0$ (GHz) & $E_{\rm u}$ & $n$ \\
\hline
SO & $3_{3}$--$3_{2}$ & 339.341459 & 26 &  \\
CN & $3_{7/2}$--$2_{5/2}$ & 340.247770 & 33 & 5 \\
SO & $7_{8}$--$6_{7}$ & 340.714155 & 81 &  \\
CS & $7$--$6$ & 342.882850 & 66 &  \\
SO & $8_{8}$--$7_{7}$ & 344.310612 & 87 &  \\
SO & $2_{3}$--$2_{1}$ & 345.704474 & 21 &  \\
CO & $3$--$2$ & 345.795990 & 33 &  \\
SO & $9_{8}$--$8_{7}$ & 346.528481 & 79 &  \\
SiO & $8$--$7$ & 347.330581 & 75 &  \\
HCN & $4$--$3$ & 354.505478 & 43 & 7 \\
HCO$^{+}$ & $4$--$3$ & 356.734223 & 43 &  \\
H$_2$CO & $7_{1,7}$--$6_{1,6}$ & 491.968369 & 106 &  \\
{[C\,\textsc{i}]} & $^3P_1$--$^3P_0$ & 492.160651 & 24 &  \\
CN & $6_{11/2}$--$5_{11/2}$ & 678.813471 & 114 & 5 \\
CN & $6_{11/2}$--$5_{9/2}$ & 680.046846 & 114 & 6 \\
CN & $6_{13/2}$--$5_{11/2}$ & 680.263967 & 114 & 6 \\
SO & $9_{10}$--$8_{8}$ & 682.684375 & 120 &  \\
\hline
\end{tabular}%

\end{table*}

\emph{Two coincidences need saying plainly.} Because the list is complete by
rule it holds transitions too weak to be plausible emitters in a debris disc,
and \FrWeakOnlyN{} crossing (\FrWeakOnlyStar) coincides with nothing stronger;
nothing turns on that, since the label does not dispose. The second is
\FrCpStar, which lies $\FrCpSODv$\,km\,s$^{-1}$ from \FrSOName, well inside the
mask, and is nevertheless reported as unattributed. The attribution is not
astrophysically credible: sulphur monoxide has not been detected in a debris
disc, where the second-generation gas is CO, C and O, and an \RcCpFlux\,mJy SO
line would here be unaccompanied by carbon monoxide. CO was covered. The same
Band~\FrCpCoBand{} window that holds the crossing holds CO($3\to2$), which
falls at \FrCpCoGHz\,GHz in that block and is absent to \FrCpCoLimMJy\,mJy at
the stellar position in a \FrCpCoChanKms\,km\,s$^{-1}$ channel --- the same
limit, in the same data, as the feature itself. The coincidence is also
inadmissible procedurally, because sulphur monoxide entered the frozen list
after this crossing had been examined. Reporting the crossing as unattributed
keeps a crossing unattributed where the rule would have attributed it,
which is the conservative
direction and the one consistent with fixing the criteria in advance; the
repeat block excludes a persistent carrier there in any case
(Appendix~\ref{app:radius}).

\emph{A wider list is not a better one.} Every Splatalogue transition over the
searched intervals under the same energy cut and with an Einstein coefficient
above $10^{-5}$\,s$^{-1}$ --- \MaskDbTrans{} transitions of \MaskDbCarriers{}
species --- masks \MaskDbPctUnion{} per cent of the union and moves
\FrFullNFlip{} of the \NCross{} crossings across the line,
\FrFullNFlipRankWord{} of them a crossing that outranks every control, at
$\FrFullFlipRankDv$\,km\,s$^{-1}$ from a \FrFullFlipRankChem{} transition. At
these frequencies most frequencies lie within $\pm\MaskHalfKms$\,km\,s$^{-1}$
of something in such a list, which is why the species are restricted to what
these fields are observed to emit, and why the outcome is a label.
``Coincident with a known molecular line'' is a statement about a catalogue
rather than a property of a frequency.
